%
% Polynome
%
\subsection{Polynome}
Als Illustration der starken Einschränkungen, die Holomorphie einer
komplexen Funktion auferlegt, zeigen wir in diesem Abschnitt,
dass eine holomorphe Funktion automatisch ein Polynom in der
komplexen Variable $z=x+iy$ ist,
wenn der Real- oder Imaginärteil ein Polynom von $x$ und $y$ ist.

\begin{satz}
\label{buch:holomorph:holomorph:satz:polynom}
Sei $f(z)$ eine holomorphe Funktion, deren Realteil ein Polynom
im Realteil $x$ und dem Imaginärteil $y$ des Arguments $z=x+iy$ ist.
Dann ist $f(z)$ ein Polynom.
\end{satz}

Um den Satz~\ref{buch:holomorph:holomorph:satz:polynom}
zu beweisen verwenden wir zuerst die Tatsache, dass der Realteil
eine harmonische Funktion ist, um zu zeigen, sich die Terme gleichen
Grades in $x$ und $y$ als Realteil der Potenzen von $z=x+iy$ schreiben
lassen.
Damit ist dann wegen Satz~\ref{buch:holomorph:holomorph:satz:utov}
auch der Imaginärteil bis auf eine Konstante bestimmt.
Der Imaginärteil des gefundenen Polynoms kann sich daher höchstens
durch eine Konstante vom Imaginärteil von $f(z)$ unterscheiden.
Dies zeigt, dass $f(z)$ ein Polynom ist.

\begin{proof}
Nach Voraussetzung ist $u(x,y)$ ein Polynom, welches wir in der Form
\[
u(x,y)
=
\sum_{k,l}
a_{kl} x^ky^l
\]
schreiben, wobei nur endlich viele der $a_{kl}$ von $0$ verschieden
sind.
Der Grad des Terms $a_{kl}x^ky^l$ ist $k+l$. 

Die Funktion $u(x,y)$ ist als Realteil einer holomorphen Funktion
harmonisch, d.~h.~$\Delta u=0$.
Für die Terme vom Grad $0$ und $1$ gilt dies automatisch, da die
zweiten partiellen Ableitungen verschwinden.
Die höchstens linearen Terme haben daher die Form
\begin{align*}
a_{00} + a_{10}x + a_{01}y
&=
a_{00} + \operatorname{Re} (a_{10}-ia_{01})(x+iy)
=
a_{00} + \operatorname{Re} a_1z
\end{align*}
mit $a_1=a_{10}-ia_{01}$.
Die höchstens linearen Terme sind also tatsächlich der Realteil eines
Polynoms ersten Grades in der Variablen $z$.

Für die Terme vom Grad $2$ verschwinden die zweiten partiellen 
Ableitungen jedoch nicht automatisch, vielmehr gilt
\begin{align*}
0
&=
\Delta
\sum_{k+l=2} a_{kl}x^ky^l
\\
&=
\Delta (a_{20}x^2 + a_{11}xy + a_{02}y^2)
\\
&=
\frac{\partial^2}{\partial x^2}
a_{20}x^2
+
\frac{\partial^2}{\partial y^2}
a_{02}y^2
\\
0
&=
2\cdot 1\cdot a_{20}  + 2 \cdot 1 \cdot a_{02}
\qquad\Rightarrow\qquad a_{02}=-a_{20}.
\end{align*}
Dies zeigt, dass die beiden Koeffizienten $a_{20}$ und $a_{02}$ nicht
unabhängig gewählt werden können, sondern dass die Terme zweiten
Grades nur in der Kombination
\(
a_{20}(x^2-y^2) + a_{11}xy
\)
vorkommen können.
Wegen
\(
z^2
=
(x^2-y^2) + 2ixy
\)
lässt sich dies auch als der Realteil 
\[
a_{20}(x^2-y^2) + a_{11}xy
=
\operatorname{Re} z^2\cdot \biggl(a_{20}-\frac{i}2a_{11}\biggr)
=
\operatorname{Re} a_2z^2
\qquad\text{mit $a_2=a_{20}-\displaystyle\frac{i}2a_{11}$}
\]
verstehen.
Auch für die Terme vom Grad 2 hat sich ergeben, dass sie der Realteil
eines komplexen Monoms in der Variablen $z$ vom Grad $2$ sind.

Das beobachtete Muster gilt jedoch allgemeiner.
Wir zeigen jetzt, dass die Terme vom Grad $n$ in $u(x,y)$ zusammen
den Realteil eines Monoms der Form $a_nz^n$ bilden.
Durch Subtraktion dieses Terms bleiben dann nur noch Terme vom Grad $n-1$,
auf die die gleiche Beobachtung wieder angewendet werden kann.
Auf diese Weise lässt sich schrittweise ein Polynom
\[
p(z)
=
a_{00} + a_1z + a_2z^2 + \dots + a_{n-1}z^{n-1} + a_nz^n
\]
mit
der Eigenschaft $u(x,y)=\operatorname{Re}p(z)$ konstruieren.
Der Imaginärteil von $f(z)$ ist durch den Realteil bis auf eine imaginäre
Konstante festgelegt ist, die wir mit $ib$ bezeichnen können.
Das Polynom
\[
q(z)
=
p(z)+ib
=
\underbrace{a_{00}+ib}_{\displaystyle a_0}
+ a_1z + a_2z^2 + \dots + a_{n-1}z^{n-1} + a_nz^n
=
\sum_{m=0}^n a_mz^m
\]
ist dann das gesuchte Polynom mit $f(z)=q(z)$.

Wir betrachten jetzt ausschliesslich die Terme
\input{chapters/020-holomorph/fig/fig-polynome.tex}
\[
u^n(x,y)
=
\sum_{k+l=n} a_{kl}x^ky^l
\]
vom Grad $n>2$ in $u(x,y)$.
Der Laplace-Operator reduziert den Grad jeweils um $2$, die entstehenden
Terme müssen sich aber alle wegheben.
Dies führt auf die Bedingung
\[
\Delta
u^n(x,y)
=
\sum_{k+l=n}
\bigl(k(k-1)a_{kl}x^{k-2}y^l
+
l(l-1)a_{kl}x^ky^{l-2}
\bigr)
=
0.
\]
In Abbildung~\ref{buch:holomorph:holomorph:fig:polynome}
\input{chapters/020-holomorph/fig/fig-rekursion.tex}
werden die einzelnen Monome als Gitterpunkte dargestellt, die Pfeile
beschreiben die Wirkung der zweiten partiellen Ableitungen.
Dort wo sich zwei Pfeilspitzen treffen, müssen sich Terme wegheben, 
wie dies in Abbildung~\ref{buch:holomorph:holomorph:fig:rekursion}
illustriert ist.
Dies geschieht jeweils für Terme, deren Exponenten sich um 2 unterscheiden:
\begin{align}
0
&=
\frac{\partial^2}{\partial x^2}
a_{k,n-k}x^ky^{n-k}
+
\frac{\partial^2}{\partial y^2}
a_{k-2,n-k+2}x^{k-2}y^{n-k+2}.
\notag
\\
&=
\bigl(
k(k-1) a_{k,n-k} + (n-k+2)(n-k+1) a_{k-2,n-k+2}
\bigr)
x^{k-2}y^{n-k}.
\notag
\intertext{Wir lesen daraus die Rekursionsformeln}
a_{k-2,n-k+2}
&=
-
\frac{
k(k-1)
}{
(n-k+2)(n-k+1)
}
a_{k,n-k}
\label{buch:holomorph:holomorph:beweis:polynom:rek1}
\\
a_{k,n-k}
&=
-
\frac{
(n-k+2)(n-k+1)
}{
k(k-1)
}
a_{k-2,n-k+2}
\label{buch:holomorph:holomorph:beweis:polynom:rek2}
\end{align}
für die Koeffizienten $a_{kl}$ ab.
Die Koeffizienten sind somit nicht unabhängig voneinander, sondern werden
durch die Rekursionsformeln
\eqref{buch:holomorph:holomorph:beweis:polynom:rek1}
und
\eqref{buch:holomorph:holomorph:beweis:polynom:rek2}
miteinandern verknüpft.

Für ungerades $n$ lassen sich alle Koeffizienten durch $a_{n0}$ oder $a_{0n}$
ausdrücken.
Es gilt nämlich nach wiederholter Anwendung von
\eqref{buch:holomorph:holomorph:beweis:polynom:rek1}
der Reihe nach
\begin{align*}
a_{n-2,2}
&=
-
\frac{n(n-1)}{2\cdot 1} a_{n0}
=
-\binom{n}{2}
\\
a_{n-4,4}
&=
-
\frac{(n-2)(n-3)}{4\cdot 3}
a_{n-2,2}
=
\frac{(n-2)(n-3)}{4\cdot 3}
\frac{n(n-1)}{2\cdot 1}
a_{n0}
=
\binom{n}{4}
a_{n0}
\\
&\phantom{0}\vdots
\\
a_{n-2k,2k}
&=
(-1)^k
\binom{n}{2k} a_{n0}
\intertext{bzw.~durch Anwendung von
\eqref{buch:holomorph:holomorph:beweis:polynom:rek2}}
a_{2,n-2}
&=
-
\frac{n(n-1)}{2\cdot 1}
a_{0n}
=
-\binom{n}{2} a_{0n}
\\
a_{4,n-4}
&=
-
\frac{(n-2)(n-3)}{4\cdot 3}
a_{2,n-2}
=
\frac{(n-2)(n-3)}{4\cdot 3}
\frac{n(n-1)}{2\cdot 1}
a_{2,n-2}
=
\binom{n}{4}a_{0n}
\\
&\phantom{0}\vdots
\\
a_{2k,n-2k}
&=
(-1)^k\binom{n}{2k} 
a_{0n}.
\end{align*}
Damit kann man jetzt $u^n(x,y)$
als
\begin{align*}
u^n(x,y)
&=
a_{n0}
\sum_{k=0}^{\lfloor\frac{n}2\rfloor}
(-1)^k
\binom{n}{2k}
x^{n-2k}y^{2k}
+
a_{0n}
\sum_{k=0}^{\lfloor\frac{n}2\rfloor}
(-1)^k
\binom{n}{2k}
x^{2k}y^{n-2k}
\intertext{schreiben.
Die erste Summe enthält die geraden Potenzen von $y$, die zweite die
ungeraden.
Man kann sie auch durch Aufteilung der Summe der binomischen Formel als}
(x+iy)^n
&=
\sum_{l=0}^n i^{n-l}\binom{n}{l} x^k y^{n-l}
\\
&=
\sum_{k=0}^{\lfloor\frac{n}2\rfloor}
(-1)^k
\binom{n}{2k}
x^{n-2k}y^{2k}
+
i^n
\sum_{k=0}^{\lfloor\frac{n}2\rfloor}
(-1)^k
\binom{n}{2k}
x^{2k}y^{n-2k}
\end{align*}
erhalten.
Somit folgt
\[
u^n(x,y)
=
\operatorname{Re}\Bigl(
(\underbrace{a_{n0}-i^n a_{0n}}_{\displaystyle=a_n})
(x+iy)^n
\Bigr)
=
\operatorname{Re} a_n
(x+iy)^n,
\]
der Term $u^n(x,y)$ ist also Realteil eines Monoms.

Für gerade $n=2m$ lassen sich ähnlich alle Terme durch $a_{n0}$ und
$a_{n-1,1}$ ausdrücken.
Wie im Falle ungeraden $n$s folgt
\[
a_{n-2k,2k}
=
(-1)^k
\binom{n}{2k}
a_{n0}.
\]
Für die verbleibenden Koeffizienten gilt durch wiederholte Anwendung
von \eqref{buch:holomorph:holomorph:beweis:polynom:rek1}
der Reihe nach
\begin{align*}
a_{n-3,3}
&=
-
\frac{
(n-1)(n-2)
}{
3\cdot 2
}
a_{n-1,1}
=
-
\binom{n}{3} \frac{1}{n} a_{n-1,1}
\\
a_{n-5,5}
&=
-
\frac{
(n-3)(n-4)
}{
5\cdot 4
}
a_{n-3,3}
=
\frac{
(n-3)(n-4)
}{
5\cdot 4
}
\frac{
(n-1)(n-2)
}{
3\cdot 2
}
a_{n-1,1}
=
\binom{n}{5}\frac{1}{n}a_{n-1,1}
\\
&\phantom{0}\vdots
\\
a_{n-2k-1,2k+1}
&=
(-1)^k
\binom{n}{2k+1}\frac{1}{n}a_{n-1,1}.
\end{align*}
Auch in diesem Fall lässt sich durch die Aufteilung
\begin{align*}
(x+iy)^n
&=
\sum_{l=0}^n i^{n-l} \binom{n}{l} x^ly^{n-l}
\\
&=
\sum_{k=0}^{\frac{n}2} (-1)^k\binom{n}{2k} x^{2k}y^{n-2k}
+
i^{n-1}
\sum_{k=0}^{\frac{n}2} (-1)^k\binom{n}{2k+1} x^{n-2k-1}y^{2k+1}
\end{align*}
in Real- und Imaginärteil der Potenz $z^n$ erkennen, dass 
$u^n(x,y)$ wieder als Realteil einer Potenz $a_nz^n$ geschrieben
werden kann.
Wie oben ausgeführt, kann man daher schliessen, dass $f(z)$ ein
Polynom ist, welches durch $u(x,y)$ bis auf eine imaginäre Konstante
festgelegt ist.
\end{proof}

